\documentclass[11pt,a4paper]{article}
\usepackage[utf8]{inputenc}
\usepackage[T1]{fontenc}
\usepackage[french]{babel}
\usepackage[margin=1.8cm]{geometry}
\usepackage{amsmath,amssymb}
\usepackage{xcolor}
\usepackage{colortbl}
\usepackage{tcolorbox}
\usepackage{booktabs}
\usepackage{array}
\usepackage{multicol}
\usepackage{enumitem}
\usepackage{tikz}
\usepackage{pgfplots}
\pgfplotsset{compat=1.17}
\usepackage{fancyhdr}
\usepackage{multirow}
\usepackage{fontawesome5}
\usepackage{longtable}

\tcbuselibrary{skins,breakable}

% ---- Couleurs ----
\definecolor{biocolor}{RGB}{0,110,110}
\definecolor{bioaccent}{RGB}{200,80,40}
\definecolor{biolight}{RGB}{230,245,245}
\definecolor{bioyellow}{RGB}{255,245,210}

% ---- En-têtes ----
\pagestyle{fancy}
\fancyhf{}
\fancyhead[L]{\textcolor{biocolor}{\textbf{Challenge Statistique -- STS BIOAL}}}
\fancyhead[R]{\textcolor{biocolor}{Module 5 -- Statistique descriptive}}
\fancyfoot[C]{\thepage}
\renewcommand{\headrulewidth}{0.6pt}
\renewcommand{\headrule}{\hbox to\headwidth{\color{biocolor}\leaders\hrule height \headrulewidth\hfill}}

% ---- Boîtes ----
\newtcolorbox{epreuve}[2][]{
  breakable, colback=biolight, colframe=biocolor,
  coltitle=white, title={\textbf{#2}},
  fonttitle=\bfseries\large, boxrule=0.8pt, arc=3pt,
  left=6pt, right=6pt, top=6pt, bottom=6pt, #1
}

\newtcolorbox{regle}{
  colback=bioyellow, colframe=bioaccent,
  boxrule=0.8pt, arc=3pt,
  left=6pt, right=6pt, top=6pt, bottom=6pt
}

\newtcolorbox{corrige}[1][]{
  breakable, colback=white, colframe=bioaccent,
  title={\textbf{Corrigé indicatif}}, coltitle=white, fonttitle=\bfseries,
  boxrule=0.8pt, arc=3pt, #1
}

% ---- Vignettes ----
\newtcolorbox{vignetteCours}[1]{
  breakable, enhanced, colback=biolight, colframe=biocolor,
  coltitle=white, fonttitle=\bfseries, title={\faBookOpen\ \ #1},
  boxrule=0.7pt, arc=4pt, left=6pt, right=6pt, top=5pt, bottom=5pt,
  drop fuzzy shadow=black!15, width=\linewidth
}

\newtcolorbox{vignetteMethode}[1]{
  breakable, enhanced, colback=bioyellow, colframe=bioaccent,
  coltitle=white, fonttitle=\bfseries, title={\faTools\ \ #1},
  boxrule=0.7pt, arc=4pt, left=6pt, right=6pt, top=5pt, bottom=5pt,
  drop fuzzy shadow=black!15, width=\linewidth
}

\newtcolorbox{vignetteAlerte}[1]{
  breakable, enhanced, colback=red!5, colframe=red!70!black,
  coltitle=white, fonttitle=\bfseries, title={\faExclamationTriangle\ \ #1},
  boxrule=0.7pt, arc=4pt, left=6pt, right=6pt, top=5pt, bottom=5pt,
  drop fuzzy shadow=black!15, width=\linewidth
}

\newtcolorbox{vignetteFormule}{
  breakable, enhanced, colback=white, colframe=biocolor!70,
  boxrule=0.6pt, arc=3pt, left=5pt, right=5pt, top=4pt, bottom=4pt,
  fontupper=\small
}

% ---- Données ----
\newcommand{\dataS}{%
12{,}1 ; 12{,}6 ; 13{,}0 ; 13{,}3 ; 13{,}6 ; 13{,}9 ; 14{,}0 ; 14{,}2 ;
14{,}4 ; 14{,}5 ; 14{,}7 ; 14{,}9 ; 15{,}0 ; 15{,}1 ; 15{,}2 ; 15{,}3 ;
15{,}4 ; 15{,}5 ; 15{,}6 ; 15{,}8 ; 15{,}9 ; 16{,}0 ; 16{,}1 ; 16{,}3 ;
16{,}4 ; 16{,}5 ; 16{,}6 ; 16{,}8 ; 16{,}9 ; 17{,}0 ; 17{,}2 ; 17{,}4 ;
17{,}5 ; 17{,}7 ; 17{,}9 ; 18{,}1 ; 18{,}4 ; 18{,}6 ; 18{,}9 ; 19{,}1}

% =========================================================
\begin{document}

% ============== PAGE DE GARDE ==============
\begin{center}
{\color{biocolor}\rule{\linewidth}{1.2pt}}\\[0.3cm]
{\Huge\bfseries\color{biocolor} CHALLENGE STATISTIQUE}\\[0.2cm]
{\Large\bfseries\color{bioaccent} STS BIOAL -- Module 5}\\[0.1cm]
{\large Statistique descriptive}\\[0.2cm]
{\color{biocolor}\rule{\linewidth}{1.2pt}}
\end{center}

\vspace{0.3cm}

\begin{regle}
\textbf{\large Règles du jeu}\\[0.2cm]
\begin{itemize}[leftmargin=*,itemsep=2pt]
  \item \textbf{Format} : équipes de \textbf{3 ou 4 étudiants}.
  \item \textbf{Durée totale} : \textbf{50 minutes} (chronométrées, épreuve par épreuve).
  \item \textbf{6 épreuves} -- \textbf{100 points} au total.
  \item \textbf{Matériel autorisé} : calculatrice, tableur (ordinateur), papier millimétré, règle.
  \item \textbf{Barème} : les points sont attribués par épreuve selon les critères indiqués. En cas d'égalité, l'équipe ayant terminé la dernière épreuve en premier gagne.
  \item Chaque équipe remet \textbf{une seule feuille-réponse} par épreuve.
\end{itemize}
\end{regle}

\vspace{0.3cm}

\begin{center}
\renewcommand{\arraystretch}{1.4}
\begin{tabular}{|c|l|c|c|}
\hline
\textbf{\'Epreuve} & \textbf{Thème} & \textbf{Durée} & \textbf{Barème} \\
\hline
1 & Effectifs, fréquences, fréquences cumulées & 8 min & 20 pts \\
2 & Diagrammes, histogramme, boîte à moustaches & 8 min & 15 pts \\
3 & Indicateurs de position & 8 min & 20 pts \\
4 & Indicateurs de dispersion & 7 min & 15 pts \\
5 & Ajustement affine et corrélation & 12 min & 20 pts \\
6 & Utilisation du tableur & 7 min & 10 pts \\
\hline
\multicolumn{2}{|r|}{\textbf{TOTAL}} & \textbf{50 min} & \textbf{100 pts} \\
\hline
\end{tabular}
\end{center}

\vspace{0.3cm}

\newpage

% =========================================================
% CONTEXTE COMMUN
% =========================================================
\begin{tcolorbox}[colback=bioyellow,colframe=bioaccent,title={\textbf{Contexte commun à toutes les épreuves}},fonttitle=\bfseries]
Un laboratoire de \textbf{bioanalyse} contrôle la \textbf{concentration en glucose} (en g/L) d'un lot de 40 échantillons sériques. Les résultats obtenus sont les suivants :

\vspace{0.2cm}
\centering
\begin{minipage}{0.95\linewidth}
\small\dataS
\end{minipage}
\end{tcolorbox}

\vspace{0.2cm}

% =========================================================
% ÉPREUVE 1
% =========================================================
\begin{epreuve}{ÉPREUVE 1 -- Tableau statistique \hfill \small (8 min -- 20 pts)}
On regroupe les 40 valeurs en \textbf{8 classes} de largeur 1 g/L :
$[12\,;13[$, $[13\,;14[$, \dots, $[19\,;20[$.

\vspace{0.2cm}
\textbf{1.1.} Compléter le tableau ci-dessous (arrondir les fréquences à $10^{-3}$) :

\vspace{0.2cm}
\renewcommand{\arraystretch}{1.5}
\begin{center}
\begin{tabular}{|c|c|c|c|c|}
\hline
\textbf{Classe} & \textbf{Effectif $n_i$} & \textbf{Fréquence $f_i$} & \textbf{Fréq. cumulée croissante} & \textbf{Centre $x_i$} \\
\hline
$[12\,;13[$ & & & & 12{,}5 \\
\hline
$[13\,;14[$ & & & & \\
\hline
$[14\,;15[$ & & & & \\
\hline
$[15\,;16[$ & & & & \\
\hline
$[16\,;17[$ & & & & \\
\hline
$[17\,;18[$ & & & & \\
\hline
$[18\,;19[$ & & & & \\
\hline
$[19\,;20[$ & & & & \\
\hline
\textbf{Total} & \textbf{40} & \textbf{1} & --- & --- \\
\hline
\end{tabular}
\end{center}

\vspace{0.2cm}
\textbf{1.2.} Quelle est la \textbf{fréquence cumulée croissante} de la classe $[16\,;17[$ ? Interpréter en une phrase.

\vspace{0.2cm}
\textbf{1.3.} Quelle proportion d'échantillons a une concentration \textbf{strictement inférieure à 15 g/L} ? Exprimer en \% (arrondi à 0{,}1\,\%).

\vspace{0.2cm}
\textbf{1.4.} Quelle proportion a une concentration \textbf{supérieure ou égale à 17 g/L} ?
\end{epreuve}

\newpage

% =========================================================
% ÉPREUVE 2
% =========================================================
\begin{epreuve}{ÉPREUVE 2 -- Représentations graphiques \hfill \small (8 min -- 15 pts)}
\textbf{2.1.} Sur le papier millimétré fourni, tracer l'\textbf{histogramme} des effectifs (classes de largeur 1, donc hauteurs = effectifs).

\vspace{0.2cm}
\textbf{2.2.} Tracer la \textbf{courbe des fréquences cumulées croissantes} sur le même graphique (utiliser les bornes supérieures des classes comme abscisses).

\vspace{0.2cm}
\textbf{2.3.} À partir du tableau de l'épreuve 1, déterminer :
\begin{itemize}[leftmargin=*,itemsep=1pt]
  \item le \textbf{minimum} : \dotfill
  \item le \textbf{maximum} : \dotfill
  \item la \textbf{médiane} $M_e$ : \dotfill
  \item le \textbf{premier quartile} $Q_1$ : \dotfill
  \item le \textbf{troisième quartile} $Q_3$ : \dotfill
\end{itemize}

\vspace{0.2cm}
\textbf{2.4.} Tracer la \textbf{boîte à moustaches} (diagramme en boîte) correspondante sur l'axe gradué ci-dessous :

\vspace{0.4cm}
\begin{center}
\begin{tikzpicture}[xscale=1.6]
  \draw[thick,<->] (11.5,0) -- (19.8,0);
  \foreach \x in {12,13,14,15,16,17,18,19}
     \draw[thick] (\x,0.1) -- (\x,-0.1) node[below]{\small \x};
  \node[left] at (11.5,0) {\small g/L};
\end{tikzpicture}
\end{center}

\vspace{0.2cm}
\textbf{2.5.} D'après la boîte à moustaches, la distribution est-elle symétrique ? Justifier.
\end{epreuve}

\newpage

% =========================================================
% ÉPREUVE 3
% =========================================================
\begin{epreuve}{ÉPREUVE 3 -- Indicateurs de position \hfill \small (8 min -- 20 pts)}
\vspace{0.1cm}
\textbf{3.1.} Calculer la \textbf{moyenne} $\bar{x}$ de la série (à partir des centres de classes). Arrondir à $10^{-2}$.

\vspace{0.2cm}
\textbf{3.2.} Déterminer la \textbf{médiane} $M_e$. Justifier brièvement la méthode utilisée (position dans la série ordonnée).

\vspace{0.2cm}
\textbf{3.3.} Déterminer $Q_1$ et $Q_3$. Rappeler la signification concrète de $Q_1$.

\vspace{0.2cm}
\textbf{3.4.} Quel est le \textbf{mode} de la série ? S'agit-il d'un mode unique ? Préciser la \textbf{classe modale}.

\vspace{0.2cm}
\textbf{3.5.} \textbf{Vrai/Faux} -- justifier :
\begin{itemize}[leftmargin=*,itemsep=2pt]
  \item \emph{« Comme la moyenne et la médiane sont proches, la distribution est approximativement symétrique. »}
  \item \emph{« 50\,\% des échantillons ont une concentration comprise entre $Q_1$ et $Q_3$. »}
\end{itemize}
\end{epreuve}

\vspace{0.3cm}

% =========================================================
% ÉPREUVE 4
% =========================================================
\begin{epreuve}{ÉPREUVE 4 -- Indicateurs de dispersion \hfill \small (7 min -- 15 pts)}
\textbf{4.1.} Calculer l'\textbf{étendue} de la série.

\vspace{0.2cm}
\textbf{4.2.} Calculer l'\textbf{écart interquartile} $I_Q = Q_3 - Q_1$.

\vspace{0.2cm}
\textbf{4.3.} Calculer la \textbf{variance} $V$ et l'\textbf{écart-type} $\sigma$ de la série (à partir des centres de classes). Arrondir $V$ à $10^{-2}$ et $\sigma$ à $10^{-3}$.

\emph{Aide : on rappelle que } $V = \dfrac{1}{N}\displaystyle\sum_{i=1}^{k} n_i\,x_i^2 - \bar{x}^{\,2}$.

\vspace{0.2cm}
\textbf{4.4.} Interpréter l'écart-type en une phrase (que mesure-t-il concrètement ici ?).
\end{epreuve}

\newpage

% =========================================================
% ÉPREUVE 5
% =========================================================
\begin{epreuve}{ÉPREUVE 5 -- Ajustement affine et corrélation \hfill \small (12 min -- 20 pts)}
Pour \textbf{étalonner} un spectrophotomètre, le laboratoire prépare 6 solutions de glucose de concentrations connues et mesure l'absorbance :

\vspace{0.2cm}
\begin{center}
\renewcommand{\arraystretch}{1.3}
\begin{tabular}{|c|c|c|c|c|c|c|}
\hline
\textbf{Concentration $x_i$ (g/L)} & 0 & 2 & 4 & 6 & 8 & 10 \\
\hline
\textbf{Absorbance $y_i$} & 0{,}05 & 0{,}22 & 0{,}38 & 0{,}57 & 0{,}71 & 0{,}90 \\
\hline
\end{tabular}
\end{center}

\vspace{0.2cm}
\textbf{5.1.} Calculer les moyennes $\bar{x}$ et $\bar{y}$.

\vspace{0.2cm}
\textbf{5.2.} Calculer les sommes suivantes :
$\displaystyle\sum x_i^2$, $\displaystyle\sum y_i^2$, $\displaystyle\sum x_i y_i$.

\vspace{0.2cm}
\textbf{5.3.} Déterminer par la \textbf{méthode des moindres carrés} l'équation de la droite de régression $y = ax + b$ (arrondir $a$ à $10^{-4}$ et $b$ à $10^{-3}$).

\vspace{0.2cm}
\textbf{5.4.} Calculer le \textbf{coefficient de corrélation linéaire} $r$ (arrondi à $10^{-3}$).

\vspace{0.2cm}
\textbf{5.5.} Le modèle linéaire est-il adapté ? Justifier à l'aide de $r$.

\vspace{0.2cm}
\textbf{5.6.} \emph{Application} : on mesure une absorbance de $0{,}80$. Estimer la concentration en glucose de l'échantillon (arrondie à $10^{-2}$ g/L).
\end{epreuve}

\vspace{0.3cm}

% =========================================================
% ÉPREUVE 6
% =========================================================
\begin{epreuve}{ÉPREUVE 6 -- Utilisation du tableur \hfill \small (7 min -- 10 pts)}
Sur la feuille de calcul, les concentrations sont saisies dans la plage \texttt{A2:A41}.

\vspace{0.2cm}
Écrire la \textbf{formule exacte} à saisir pour obtenir :

\vspace{0.2cm}
\begin{center}
\renewcommand{\arraystretch}{1.5}
\begin{tabular}{|l|l|}
\hline
\textbf{Question} & \textbf{Formule} \\
\hline
\textbf{a.} L'effectif total & \\
\hline
\textbf{b.} La moyenne & \\
\hline
\textbf{c.} La médiane & \\
\hline
\textbf{d.} Le premier quartile $Q_1$ & \\
\hline
\textbf{e.} Le troisième quartile $Q_3$ & \\
\hline
\textbf{f.} L'écart-type & \\
\hline
\textbf{g.} Le minimum & \\
\hline
\textbf{h.} Le maximum & \\
\hline
\end{tabular}
\end{center}

\vspace{0.2cm}
\textbf{i.} Citer deux méthodes pour tracer une \textbf{boîte à moustaches} à partir d'une série de données.

\vspace{0.2cm}
\textbf{j.} Après avoir entré la formule \texttt{=MOYENNE(A2:A41)} dans la cellule \texttt{C1}, on la recopie vers le bas en \texttt{C2}. Quelle formule obtient-on ? Le résultat est-il pertinent ? Que faire ?
\end{epreuve}

\newpage

% =========================================================
% FEUILLE DE SCORE (mise en page améliorée)
% =========================================================
\section*{\color{biocolor}Feuille de score de l'équipe}

% ---------- Bloc 1 : identification ----------
\begin{tcolorbox}[enhanced, colback=biolight, colframe=biocolor,
                  title={\faUsers\ \ \textbf{Identification de l'équipe}},
                  coltitle=white, fonttitle=\bfseries\large,
                  boxrule=0.8pt, arc=4pt, left=10pt, right=10pt, top=8pt, bottom=8pt,
                  drop fuzzy shadow=black!15]
\large
\textbf{Nom de l'équipe} : \hrulefill \\[0.45cm]
\textbf{Capitaine} : \hrulefill \\[0.45cm]
\textbf{Membres} : \\[0.15cm]
\begin{enumerate}[leftmargin=1cm, itemsep=0.6cm, label=\arabic*., font=\large]
\item \hrulefill
\item \hrulefill
\item \hrulefill
\item \hrulefill
\end{enumerate}
\end{tcolorbox}

\vspace{0.4cm}

% ---------- Bloc 2 : grille des scores ----------
\begin{tcolorbox}[enhanced, colback=bioyellow, colframe=bioaccent,
                  title={\faTrophy\ \ \textbf{Grille des scores}},
                  coltitle=white, fonttitle=\bfseries\large,
                  boxrule=0.8pt, arc=4pt, left=10pt, right=10pt, top=8pt, bottom=8pt,
                  drop fuzzy shadow=black!15]
\renewcommand{\arraystretch}{1.55}
\begin{tabular}{|>{\raggedright\arraybackslash}p{6.3cm}|c|>{\centering\arraybackslash}p{2.3cm}|>{\raggedright\arraybackslash}p{3.6cm}|}
\hline
\rowcolor{bioaccent!20}
\textbf{\'Epreuve} & \textbf{Barème} & \textbf{Score} & \textbf{Observations} \\
\hline
1 -- Tableau statistique & /20 & & \\ \hline
2 -- Représentations graphiques & /15 & & \\ \hline
3 -- Indicateurs de position & /20 & & \\ \hline
4 -- Indicateurs de dispersion & /15 & & \\ \hline
5 -- Ajustement affine et corrélation & /20 & & \\ \hline
6 -- Utilisation du tableur & /10 & & \\ \hline
\rowcolor{bioaccent!25}
\multicolumn{1}{|r|}{\textbf{TOTAL}} & \textbf{/100} & & \\ \hline
\end{tabular}
\end{tcolorbox}

\vspace{0.4cm}

% ---------- Bloc 3 : citations ----------
\begin{tcolorbox}[colback=bioyellow!50,colframe=bioaccent,title={\textbf{Deux citations pour méditer sur la statistique}},fonttitle=\bfseries]
\small
\textbf{1. La formule anglaise de Ronald Coase} (économiste, prix Nobel 1991) :

\medskip
\emph{« If you torture the data long enough, it will confess. »}\\
\hfill {\itshape --- Ronald Coase, années 1960}

\medskip
\textbf{Traduction française :}\\
\emph{« Si vous torturez les données assez longtemps, elles avoueront. »}

\medskip
Première trace écrite : en 1971, dans une allocution du mathématicien britannique \textbf{Irving John Good}, qui l'attribue explicitement à Coase.

\vspace{0.3cm}
\noindent\rule{\linewidth}{0.4pt}

\vspace{0.2cm}
\textbf{2. La variante française antérieure d'Alfred Sauvy} (démographe, 1898--1990) :

\medskip
\emph{« Les chiffres sont des innocents qui, sous la sollicitation, sous la torture, avouent très vite ce qu'on leur demande, quitte à se rétracter plus tard. »}\\
\hfill {\itshape --- Alfred Sauvy, \emph{L'Opinion publique}, 1961}

\vspace{0.3cm}
\noindent\rule{\linewidth}{0.4pt}

\vspace{0.2cm}
\textbf{En résumé.} Toutes deux rappellent que les chiffres peuvent être manipulés et qu'il faut se méfier des conclusions hâtives.
\end{tcolorbox}

\newpage

% =========================================================
% ANNEXE A – VIGNETTES DE COURS
% =========================================================
\section*{\color{biocolor}Annexe A -- Vignettes de cours}
\emph{À consulter à tout moment pendant l'épreuve.}

\vspace{0.3cm}

\begin{multicols}{2}

\begin{vignetteCours}{V1 -- Vocabulaire de base}
\textbf{Population} : ensemble étudié.\\
\textbf{Individu} : élément de la population.\\
\textbf{Caractère} : ce que l'on mesure (quantitatif ou qualitatif).\\
\textbf{Série statistique} : liste des valeurs observées.

\medskip
\textbf{Effectif} $n_i$ : nombre d'individus de la modalité $i$.\\
\textbf{Effectif total} : $N = \sum n_i$.\\
\textbf{Fréquence} : $f_i = \dfrac{n_i}{N}$, avec $\sum f_i = 1$.

\medskip
\textbf{Fréquence cumulée croissante} (FCC) : $F_i = f_1+f_2+\dots+f_i$.\\
La dernière FCC vaut toujours $1$.
\end{vignetteCours}

\begin{vignetteCours}{V2 -- Classes et centres de classes}
Pour une série \textbf{continue} (ou très nombreuse), on regroupe les valeurs en \textbf{classes} $[a\,;b[$.

\medskip
\textbf{Centre de classe} : $x_i = \dfrac{a+b}{2}$.

\medskip
\textbf{Amplitude} : $b-a$ (constante si classes régulières).

\medskip
On remplace chaque valeur par le centre de sa classe pour les calculs.
\end{vignetteCours}

\begin{vignetteCours}{V3 -- Représentations graphiques}
\textbf{Diagramme en bâtons} : effectifs ou fréquences (série discrète).

\medskip
\textbf{Histogramme} : classes en abscisse, \emph{aires proportionnelles} aux effectifs.\\
Classes de même amplitude $\Rightarrow$ hauteur $=$ effectif.

\medskip
\textbf{Courbe des FCC} : en abscisse les bornes supérieures, en ordonnée les fréquences cumulées. Toujours croissante, de 0 à 1.

\medskip
\textbf{Boîte à moustaches} : représente $\min,\ Q_1,\ M_e,\ Q_3,\ \max$.
\end{vignetteCours}

\begin{vignetteCours}{V4 -- Indicateurs de position}
\textbf{Moyenne} : $\bar{x} = \dfrac{1}{N}\displaystyle\sum_{i=1}^{k} n_i x_i$.

\medskip
\textbf{Médiane} $M_e$ : valeur qui partage la série ordonnée en deux moitiés (50\,\% / 50\,\%).

\medskip
\textbf{Quartiles} : $Q_1$ (25\,\%), $M_e$ (50\,\%), $Q_3$ (75\,\%).

\medskip
\textbf{Mode} : valeur (ou classe) de plus grand effectif.
\end{vignetteCours}

\begin{vignetteCours}{V5 -- Indicateurs de dispersion}
\textbf{Étendue} : $\max - \min$.

\medskip
\textbf{Écart interquartile} : $I_Q = Q_3 - Q_1$.

\medskip
\textbf{Variance} : 
$V = \dfrac{1}{N}\displaystyle\sum n_i(x_i-\bar{x})^2
= \dfrac{1}{N}\sum n_i x_i^2 - \bar{x}^2$.

\medskip
\textbf{Écart-type} : $\sigma = \sqrt{V}$ (même unité que $x$).
\end{vignetteCours}

\begin{vignetteCours}{V6 -- Ajustement affine et corrélation}
Nuage de points $(x_i\,;y_i)$.\\
\textbf{Droite de régression} : $y = a\,x + b$ (moindres carrés).

\medskip
\textbf{Covariance} :
$\mathrm{Cov}(x,y)=\dfrac{1}{N}\sum x_i y_i - \bar{x}\,\bar{y}$.

\medskip
\textbf{Coefficient de corrélation linéaire} :
$r = \dfrac{\mathrm{Cov}(x,y)}{\sigma_x\,\sigma_y}$.

\medskip
Propriété : $-1 \le r \le 1$.\\
$|r|$ proche de $1$ $\Rightarrow$ forte corrélation linéaire.
\end{vignetteCours}

\end{multicols}

\newpage

% ---------- V7 (hors multicols pour une meilleure lisibilité) ----------
\begin{vignetteCours}{V7 -- Choisir entre (moyenne ; écart-type) et (médiane ; écart interquartile)}
\textbf{Deux couples d'indicateurs pour résumer une série quantitative :}
\begin{itemize}[leftmargin=*,itemsep=1pt]
  \item \textbf{Couple paramétrique} : moyenne $\bar{x}$ et écart-type $\sigma$ (ou variance $V$).
  \item \textbf{Couple robuste} : médiane $M_e$ et écart interquartile $I_Q = Q_3 - Q_1$.
\end{itemize}

\medskip
\textbf{Rôle spécifique de chaque couple.}
\begin{itemize}[leftmargin=*,itemsep=1pt]
  \item \textbf{Moyenne $\bar{x}$} : centre de gravité de la série ; utilise toutes les valeurs. Elle indique la valeur moyenne.
  \item \textbf{Écart-type $\sigma$} : dispersion moyenne autour de $\bar{x}$ ; même unité que la variable. Plus $\sigma$ est grand, plus les valeurs sont hétérogènes.
  \item \textbf{Médiane $M_e$} : valeur qui partage la série ordonnée en deux moitiés (50\,\% en dessous, 50\,\% au-dessus). Elle décrit une valeur typique, peu sensible aux extrêmes.
  \item \textbf{Écart interquartile $I_Q$} : étendue de la moitié centrale des données (entre $Q_1$ et $Q_3$). Il mesure la dispersion robuste, sans être influencé par les valeurs extrêmes.
\end{itemize}

\medskip
\textbf{Différences et similarités.}
\begin{center}
\small
\renewcommand{\arraystretch}{1.25}
\begin{tabular}{|p{3.2cm}|p{5.4cm}|p{5.4cm}|}
\hline
\textbf{Critère} & \textbf{Couple (moyenne ; écart-type)} & \textbf{Couple (médiane ; interquartile)} \\
\hline
Nature & Paramétrique, algébrique & Robuste, non paramétrique \\
\hline
Sensibilité aux extrêmes & Très sensible & Peu sensible \\
\hline
Information utilisée & Toutes les valeurs & Rangs centraux (50\,\% central) \\
\hline
Distribution idéale & Symétrique, sans outlier, unimodale & Asymétrique ou avec outliers \\
\hline
Position & Centre de gravité $\bar{x}$ & Valeur de rang 50\,\% $M_e$ \\
\hline
Dispersion & Écart moyen à la moyenne $\sigma$ & Étendue des 50\,\% centraux $I_Q$ \\
\hline
Interprétation & $\bar{x} \pm \sigma$ (si symétrie) & $M_e\ [Q_1\,;Q_3]$ \\
\hline
Usage typique & Calculs, régression, tests paramétriques & Description robuste, tests non paramétriques \\
\hline
Similarité & \multicolumn{2}{p{10.8cm}|}{Les deux couples résument position et dispersion, dans l'unité de la variable. Si la distribution est symétrique et sans outlier, $\bar{x} \approx M_e$ et $\sigma \approx I_Q / 1{,}35$.} \\
\hline
\end{tabular}
\end{center}

\medskip
\textbf{Quand utiliser l'un plutôt que l'autre ?}
\begin{itemize}[leftmargin=*,itemsep=1pt]
  \item \textbf{Utiliser moyenne + écart-type} si :
    \begin{itemize}[leftmargin=*,itemsep=0pt]
      \item distribution approximativement symétrique (histogramme en cloche) ;
      \item absence de valeurs aberrantes ;
      \item besoin de calculs ultérieurs (variance, régression, tests paramétriques) ;
      \item variable quantitative continue, mesures fiables.
    \end{itemize}
  \item \textbf{Utiliser médiane + écart interquartile} si :
    \begin{itemize}[leftmargin=*,itemsep=0pt]
      \item distribution asymétrique (ex. concentrations biologiques, temps d'attente) ;
      \item présence de valeurs extrêmes ou aberrantes ;
      \item effectif faible ou variable ordinale ;
      \item on veut décrire la moitié centrale sans être trompé par les extrêmes.
    \end{itemize}
  \item \textbf{Cas où ils sont interchangeables} : lorsque la distribution est symétrique, unimodale et sans outlier (proche d'une loi normale). Alors $\bar{x} \approx M_e$ et $\sigma \approx I_Q / 1{,}35$ ; les deux couples donnent le même message. En pratique, on préfère souvent $\bar{x}$ et $\sigma$ car ils utilisent toute l'information et se prêtent mieux aux calculs. Mais ils ne sont jamais totalement interchangeables pour les tests ou en présence d'asymétrie.
\end{itemize}

\medskip
\textbf{À retenir.} Toujours tracer un histogramme ou une boîte à moustaches avant de choisir. Si $\bar{x}$ et $M_e$ sont proches et que la distribution est symétrique, les deux couples sont cohérents. Si $\bar{x}$ s'éloigne de $M_e$ ou si la boîte à moustaches montre des outliers, privilégier $M_e$ et $I_Q$.
\end{vignetteCours}

\newpage

% =========================================================
% ANNEXE B – VIGNETTES METHODES
% =========================================================
\section*{\color{bioaccent}Annexe B -- Vignettes méthodes}
\emph{Procédures pas-à-pas à appliquer pendant l'épreuve.}

\vspace{0.3cm}

\begin{multicols}{2}

\begin{vignetteMethode}{M1 -- Construire un tableau statistique}
\begin{enumerate}[leftmargin=*,itemsep=1pt]
  \item Choisir les classes (souvent de même amplitude).
  \item Compter les effectifs $n_i$ (par comptage ou tableur).
  \item Calculer chaque fréquence $f_i = n_i/N$.
  \item Calculer les FCC en cumulant les $f_i$.
  \item Calculer les centres $x_i = (a+b)/2$.
  \item Vérifier : $\sum n_i = N$ et $\sum f_i = 1$.
\end{enumerate}
\end{vignetteMethode}

\begin{vignetteMethode}{M2 -- Calculer moyenne et écart-type}
\begin{enumerate}[leftmargin=*,itemsep=1pt]
  \item Construire le tableau : $n_i$, $x_i$, $n_i x_i$, $n_i x_i^2$.
  \item Moyenne : $\bar{x} = \dfrac{\sum n_i x_i}{N}$.
  \item $V = \dfrac{\sum n_i x_i^2}{N} - \bar{x}^2$.
  \item $\sigma = \sqrt{V}$.
  \item Arrondir à la fin \emph{seulement}.
\end{enumerate}
\end{vignetteMethode}

\begin{vignetteMethode}{M3 -- Trouver médiane et quartiles}
\textbf{Cas discret} (série ordonnée de $N$ valeurs) :
\begin{itemize}[leftmargin=*,itemsep=1pt]
  \item $M_e$ : si $N$ impair, rang $\frac{N+1}{2}$ ; si $N$ pair, moyenne des rangs $\frac{N}{2}$ et $\frac{N}{2}+1$.
  \item $Q_1$ : valeur de rang $\lceil N/4 \rceil$.
  \item $Q_3$ : valeur de rang $\lceil 3N/4 \rceil$.
\end{itemize}

\medskip
\textbf{Cas regroupé en classes} : utiliser la \textbf{courbe des FCC}. On lit l'abscisse pour $y=0{,}25$, $0{,}5$, $0{,}75$.
\end{vignetteMethode}

\begin{vignetteMethode}{M4 -- Tracer une boîte à moustaches}
\begin{enumerate}[leftmargin=*,itemsep=1pt]
  \item Tracer un axe gradué avec l'unité de la variable.
  \item Placer $\min$, $Q_1$, $M_e$, $Q_3$, $\max$.
  \item Dessiner le rectangle $[Q_1\,;Q_3]$.
  \item Tracer un trait vertical à $M_e$.
  \item Relier $\min$ à $Q_1$ et $Q_3$ à $\max$.
\end{enumerate}

\medskip
\textbf{Lecture} : boîte centrée + moustaches égales $\Rightarrow$ symétrie.
\end{vignetteMethode}

\begin{vignetteMethode}{M5 -- Ajustement affine par moindres carrés}
\begin{enumerate}[leftmargin=*,itemsep=1pt]
  \item Calculer $\bar{x}$ et $\bar{y}$.
  \item Calculer $\sum x_i^2$, $\sum y_i^2$, $\sum x_i y_i$.
  \item $\mathrm{Cov}(x,y)=\dfrac{\sum x_i y_i}{N}-\bar{x}\bar{y}$.
  \item $V(x)=\dfrac{\sum x_i^2}{N}-\bar{x}^2$ ; $V(y)=\dfrac{\sum y_i^2}{N}-\bar{y}^2$.
  \item $a = \dfrac{\mathrm{Cov}(x,y)}{V(x)}$ ; $b = \bar{y} - a\bar{x}$.
  \item $r = \dfrac{\mathrm{Cov}(x,y)}{\sqrt{V(x)\,V(y)}}$.
\end{enumerate}
\end{vignetteMethode}

\begin{vignetteMethode}{M6 -- Utiliser un tableur}
\textbf{Fonctions utiles} (plage \texttt{A2:A41}) :
\begin{itemize}[leftmargin=*,itemsep=1pt]
  \item \texttt{=NB(A2:A41)} -- effectif
  \item \texttt{=MOYENNE(A2:A41)} -- moyenne
  \item \texttt{=MEDIANE(A2:A41)} -- médiane
  \item \texttt{=QUARTILE(A2:A41;1)} -- $Q_1$
  \item \texttt{=QUARTILE(A2:A41;3)} -- $Q_3$
  \item \texttt{=ECARTYPE(A2:A41)} -- écart-type
  \item \texttt{=MIN(...)} / \texttt{=MAX(...)}
\end{itemize}

\medskip
\textbf{Piège} : lors d'une recopie vers le bas, la plage se décale (\texttt{A3:A42}). Utiliser des références absolues : \texttt{\$A\$2:\$A\$41}.
\end{vignetteMethode}

\end{multicols}

\vspace{0.3cm}

% ---------- Formulaire récapitulatif ----------
\section*{\color{biocolor}Annexe C -- Formulaire récapitulatif}
\begin{vignetteFormule}
\begin{multicols}{2}
\noindent\textbf{Effectif total} : $N=\sum n_i$\\
\textbf{Fréquence} : $f_i = n_i/N$\\
\textbf{FCC} : $F_i = \sum_{j\le i} f_j$\\
\textbf{Centre de classe} : $x_i=(a+b)/2$\\
\textbf{Moyenne} : $\bar{x}=\frac{1}{N}\sum n_i x_i$\\
\textbf{Variance} : $V=\frac{1}{N}\sum n_i x_i^2-\bar{x}^2$\\
\textbf{Écart-type} : $\sigma=\sqrt{V}$\\
\textbf{Étendue} : $E=\max-\min$\\
\textbf{Écart interquartile} : $I_Q=Q_3-Q_1$\\
\textbf{Covariance} : $\mathrm{Cov}=\frac{1}{N}\sum x_i y_i-\bar{x}\bar{y}$\\
\textbf{Pente} : $a=\mathrm{Cov}(x,y)/V(x)$\\
\textbf{Ordonnée} : $b=\bar{y}-a\bar{x}$\\
\textbf{Corrélation} : $r=\mathrm{Cov}(x,y)/(\sigma_x\sigma_y)$
\end{multicols}
\end{vignetteFormule}

\newpage

% =========================================================
% ANNEXE D – VIGNETTES D'EXEMPLES DE DIAGRAMMES
% =========================================================
\section*{\color{biocolor}Annexe D -- Vignettes d'exemples de diagrammes}
\emph{Exemples illustratifs \textbf{indépendants du sujet} : ils servent à comprendre chaque type de représentation, sans donner d'éléments de réponse.}

\vspace{0.3cm}

\begin{vignetteCours}{D1 -- Diagramme en bâtons (variable discrète)}
\emph{Exemple : nombre de colonies observées sur 25 boîtes de Petri.}
\begin{center}
\begin{tikzpicture}
\begin{axis}[
  width=12cm, height=5cm,
  ybar, bar width=14pt,
  xlabel={Nombre de colonies},
  ylabel={Effectif $n_i$},
  ymin=0, ymax=10,
  xtick={0,1,2,3,4,5},
  ytick={0,2,4,6,8,10},
  axis lines=left,
  tick label style={font=\small},
  label style={font=\small},
  nodes near coords,
  every node near coord/.append style={font=\small},
  every axis plot/.append style={fill=biocolor!60, draw=biocolor},
]
\addplot coordinates {(0,2)(1,5)(2,8)(3,6)(4,3)(5,1)};
\end{axis}
\end{tikzpicture}
\end{center}
\textbf{Lecture} : chaque barre correspond à une valeur \emph{unique} de la variable.\\
\textbf{Quand l'utiliser} : variable quantitative \emph{discrète} (comptage) ou qualitative.\\
\textbf{Intérêt} : visualise immédiatement \textbf{la fréquence de chaque modalité}, permet d'identifier le \textbf{mode} d'un seul coup d'œil et de comparer rapidement les catégories entre elles. Idéal pour les petits effectifs et les variables qualitatives.
\end{vignetteCours}

\begin{vignetteCours}{D2 -- Diagramme circulaire (secteurs)}
\emph{Répartition de 20 échantillons en 5 catégories (A : 30\,\%, B : 25\,\%, C : 20\,\%, D : 15\,\%, E : 10\,\%).}
\begin{center}
\begin{tikzpicture}[scale=1.15]
  \fill[bioaccent!85] (0,0) -- (0:1.7) arc (0:108:1.7) -- cycle;
  \fill[biocolor!85] (0,0) -- (108:1.7) arc (108:198:1.7) -- cycle;
  \fill[biocolor!50] (0,0) -- (198:1.7) arc (198:270:1.7) -- cycle;
  \fill[bioaccent!45] (0,0) -- (270:1.7) arc (270:324:1.7) -- cycle;
  \fill[gray!55] (0,0) -- (324:1.7) arc (324:360:1.7) -- cycle;
  \draw[thick] (0,0) circle (1.7);
  \node at (54:1.1) {\small A 30\,\%};
  \node at (153:1.1) {\small B 25\,\%};
  \node at (234:1.1) {\small C 20\,\%};
  \node at (297:1.1) {\small D 15\,\%};
  \node at (342:1.15) {\small E 10\,\%};
\end{tikzpicture}
\end{center}
\textbf{Lecture} : l'angle de chaque secteur est proportionnel à la fréquence : $\theta_i = f_i \times 360^\circ$.\\
\textbf{Quand l'utiliser} : caractère qualitatif à peu de modalités.\\
\textbf{Intérêt} : met en évidence la \textbf{part de chaque catégorie dans le total} (composition d'une population). Très parlant pour le grand public, il permet de comparer visuellement des proportions. Attention : peu lisible au-delà de 6 ou 7 modalités.
\end{vignetteCours}

\begin{vignetteCours}{D3 -- Histogramme (classes jointives)}
\emph{Exemple : longueur de feuilles (mm) mesurée sur 50 plantules (contexte totalement différent du sujet).}
\begin{center}
\begin{tikzpicture}
\begin{axis}[
  width=13cm, height=5cm,
  ybar interval,
  xlabel={Longueur de feuilles (mm)},
  ylabel={Effectif $n_i$},
  xmin=20, xmax=70, ymin=0, ymax=20,
  xtick={20,30,40,50,60,70},
  ytick={0,4,8,12,16,20},
  axis lines=left,
  tick label style={font=\small},
  label style={font=\small},
  every axis plot/.append style={fill=biocolor!55, draw=biocolor},
]
\addplot coordinates {(20,6)(30,14)(40,18)(50,9)(60,3)(70,0)};
\end{axis}
\end{tikzpicture}
\end{center}
\textbf{Lecture} : les classes sont \emph{jointives} (pas d'espace entre les barres).\\
\textbf{Règle d'or} : les \emph{aires} des rectangles sont proportionnelles aux effectifs. Si les classes ont même amplitude, hauteur $=$ effectif ; sinon, $\text{hauteur} = n_i / \text{amplitude}_i$.\\
\textbf{Intérêt} : outil privilégié pour une \textbf{variable continue regroupée en classes}. Il révèle la \textbf{forme de la distribution} : symétrie, asymétrie, unimodalité, bimodalité, présence de valeurs aberrantes. C'est le point de départ pour choisir entre moyenne et médiane.
\end{vignetteCours}

\begin{vignetteCours}{D4 -- Courbe des fréquences cumulées croissantes (FCC)}
\emph{Même exemple que D3 : longueur de feuilles (mm), $N=50$.}
\begin{center}
\begin{tikzpicture}
\begin{axis}[
  width=13cm, height=5.5cm,
  xlabel={Longueur de feuilles (mm)},
  ylabel={Fréquence cumulée croissante},
  xmin=20, xmax=70, ymin=0, ymax=1.05,
  xtick={20,30,40,50,60,70},
  ytick={0,0.25,0.5,0.75,1},
  yticklabels={0, 0{,}25, 0{,}50, 0{,}75, 1},
  grid=both, grid style={gray!25},
  axis lines=left,
  tick label style={font=\small},
  label style={font=\small},
]
\addplot[thick, color=biocolor, mark=*, mark size=2pt] 
  coordinates {(20,0)(30,0.12)(40,0.40)(50,0.76)(60,0.94)(70,1)};
\draw[dashed, bioaccent, thick] (axis cs:20,0.25) -- (axis cs:34.6,0.25) -- (axis cs:34.6,0);
\draw[dashed, bioaccent, thick] (axis cs:20,0.50) -- (axis cs:42.8,0.50) -- (axis cs:42.8,0);
\draw[dashed, bioaccent, thick] (axis cs:20,0.75) -- (axis cs:49.7,0.75) -- (axis cs:49.7,0);
\node[font=\scriptsize, bioaccent] at (axis cs:32,0.32) {$Q_1$};
\node[font=\scriptsize, bioaccent] at (axis cs:40,0.57) {$M_e$};
\node[font=\scriptsize, bioaccent] at (axis cs:47,0.82) {$Q_3$};
\end{axis}
\end{tikzpicture}
\end{center}
\textbf{Lecture} : on lit $M_e$, $Q_1$, $Q_3$ en projetant les niveaux $0{,}50$ ; $0{,}25$ ; $0{,}75$ sur l'axe des abscisses.\\
\textbf{Propriété} : la courbe est toujours croissante, de $0$ à $1$.\\
\textbf{Intérêt} : permet de \textbf{lire directement} la médiane et les quartiles (sans calcul) et de répondre à toute question du type « quel pourcentage a une valeur inférieure à $x$ ? ». Complément naturel de l'histogramme.
\end{vignetteCours}

\begin{vignetteCours}{D5 -- Boîte à moustaches (diagramme en boîte)}
\emph{Exemple : taux d'hémoglobine (g/dL) dans une cohorte de patients (contexte distinct du sujet).}
\begin{center}
\begin{tikzpicture}[xscale=1.15]
  \draw[thick, <->] (9.5,0) -- (18,0);
  \foreach \x in {10,11,12,13,14,15,16,17}
    \draw[thick] (\x,0.08) -- (\x,-0.08) node[below,font=\small]{\x};
  \draw[thick, fill=biocolor!35] (12.5, -0.45) rectangle (15.1, 0.45);
  \draw[thick] (13.8, -0.45) -- (13.8, 0.45);
  \draw[thick] (10.2, 0) -- (12.5, 0);
  \draw[thick] (15.1, 0) -- (17.4, 0);
  \draw[thick] (10.2, -0.18) -- (10.2, 0.18);
  \draw[thick] (17.4, -0.18) -- (17.4, 0.18);
  \node[above, font=\scriptsize] at (10.2, 0.2) {min};
  \node[above, font=\scriptsize] at (12.5, 0.55) {$Q_1$};
  \node[above, font=\scriptsize] at (13.8, 0.55) {$M_e$};
  \node[above, font=\scriptsize] at (15.1, 0.55) {$Q_3$};
  \node[above, font=\scriptsize] at (17.4, 0.2) {max};
  \node[below, font=\scriptsize] at (13.8, -0.65) {g/dL};
\end{tikzpicture}
\end{center}
\textbf{Lecture} : la boîte contient 50\,\% des valeurs (entre $Q_1$ et $Q_3$). La longueur des moustaches révèle l'asymétrie.\\
\textbf{Ici} : $M_e - Q_1 = 1{,}3$ et $Q_3 - M_e = 1{,}3$ $\Rightarrow$ distribution quasi symétrique.\\
\textbf{Intérêt} : \textbf{résume en une seule figure} 5 indicateurs (min, $Q_1$, $M_e$, $Q_3$, max). Permet de comparer \textbf{plusieurs séries côte à côte} et de \textbf{repérer les outliers} (points au-delà des moustaches). Idéal pour une première analyse rapide.
\end{vignetteCours}

\begin{vignetteCours}{D6 -- Nuage de points et droite de régression (cas favorable)}
\emph{Exemple : longueur (cm) et masse fraîche (g) de 6 plantules (contexte distinct du sujet).}
\begin{center}
\begin{tikzpicture}
\begin{axis}[
  width=13cm, height=6cm,
  xlabel={Longueur $x_i$ (cm)},
  ylabel={Masse fraîche $y_i$ (g)},
  xmin=0, xmax=13, ymin=0, ymax=4.5,
  xtick={0,2,4,6,8,10,12},
  ytick={0,1,2,3,4},
  grid=both, grid style={gray!25},
  axis lines=left,
  tick label style={font=\small},
  label style={font=\small},
]
\addplot[only marks, mark=*, mark size=2.5pt, color=bioaccent] 
  coordinates {(2,0.4)(4,1.2)(6,1.7)(8,2.5)(10,3.0)(12,3.8)};
\addplot[thick, color=biocolor, domain=0:13, samples=2] {0.331*x - 0.220};
\node[font=\small, bioaccent] at (axis cs:8.5,1.1) {$y \approx 0{,}331\,x - 0{,}220$};
\node[font=\small, bioaccent] at (axis cs:8.5,0.5) {$r \approx 0{,}998$};
\end{axis}
\end{tikzpicture}
\end{center}
\textbf{Lecture} : les points sont presque alignés $\Rightarrow$ $|r|$ proche de 1. On peut estimer une valeur inconnue en lisant sur la droite.\\
\textbf{Cas favorable} : nuage allongé selon une direction rectiligne.\\
\textbf{Intérêt} : visualise \textbf{la liaison entre deux variables quantitatives}. Permet de \textbf{vérifier la linéarité} avant tout calcul de régression, de \textbf{repérer les outliers} et de \textbf{prédire} une valeur $y$ pour un $x$ donné (interpolation) ou inversement (étalonnage).
\end{vignetteCours}

\begin{vignetteAlerte}{D7 -- Piège 1 : relation non linéaire (parabole)}
\emph{Exemple volontairement piégeux : $y$ dépend fortement de $x$, mais \textbf{pas linéairement}.}
\begin{center}
\begin{tikzpicture}
\begin{axis}[
  width=13cm, height=6cm,
  xlabel={Variable $x$},
  ylabel={Variable $y$},
  xmin=-0.5, xmax=10.5, ymin=-0.5, ymax=10.5,
  xtick={0,2,4,6,8,10},
  ytick={0,2,4,6,8,10},
  grid=both, grid style={gray!25},
  axis lines=left,
  tick label style={font=\small},
  label style={font=\small},
]
\addplot[only marks, mark=*, mark size=2.5pt, color=bioaccent] 
  coordinates {(0,9.5)(1,6.8)(2,4.1)(3,2.2)(4,0.9)(5,0.4)(6,0.8)(7,2.1)(8,4.0)(9,6.9)(10,9.4)};
\addplot[thick, color=biocolor, domain=0:10, samples=2] {-0.05*x + 4.3};
\node[font=\small, bioaccent] at (axis cs:7,7) {droite inadaptée};
\node[font=\small, bioaccent] at (axis cs:7,6.2) {$r \approx 0{,}02$};
\end{axis}
\end{tikzpicture}
\end{center}
\textbf{Constat} : le nuage dessine clairement une \textbf{parabole} (relation en U). Pourtant $r \approx 0$ car la droite de régression est quasi horizontale.\\
\textbf{Leçon} : un $r$ proche de 0 \textbf{ne prouve pas} l'absence de lien entre $x$ et $y$ : il prouve seulement l'absence de lien \emph{linéaire}.\\
\textbf{Reflexe} : \textbf{toujours tracer le nuage} avant de calculer $r$.
\end{vignetteAlerte}

\begin{vignetteAlerte}{D8 -- Piège 2 : absence totale de corrélation}
\emph{Exemple : deux variables totalement indépendantes (aucune tendance).}
\begin{center}
\begin{tikzpicture}
\begin{axis}[
  width=13cm, height=6cm,
  xlabel={Variable $x$},
  ylabel={Variable $y$},
  xmin=0, xmax=11, ymin=0, ymax=10,
  xtick={0,2,4,6,8,10},
  ytick={0,2,4,6,8,10},
  grid=both, grid style={gray!25},
  axis lines=left,
  tick label style={font=\small},
  label style={font=\small},
]
\addplot[only marks, mark=*, mark size=2.5pt, color=bioaccent] 
  coordinates {(1,5.3)(2,4.2)(3,5.8)(4,4.5)(5,5.1)(6,4.8)(7,5.5)(8,4.3)(9,5.7)(10,4.6)};
\addplot[thick, color=biocolor, domain=0:10, samples=2] {-0.02*x + 5.1};
\node[font=\small, bioaccent] at (axis cs:6,7.8) {aucune tendance};
\node[font=\small, bioaccent] at (axis cs:6,7) {$r \approx -0{,}10$};
\end{axis}
\end{tikzpicture}
\end{center}
\textbf{Constat} : nuage en « nuage », sans direction privilégiée. $r$ très proche de 0.\\
\textbf{Leçon} : la droite de régression passe par le centre du nuage mais \textbf{ne résume rien}. Toute prévision à partir de cette droite serait \textbf{dépourvue de sens}.\\
\textbf{Reflexe} : si $|r| < 0{,}5$, ne pas utiliser de modèle linéaire pour prévoir.
\end{vignetteAlerte}

\begin{vignetteAlerte}{D9 -- Piège 3 : point aberrant}
\emph{Exemple : 7 points parfaitement alignés sur $y = 2x$, + \textbf{1 point aberrant}.}
\begin{center}
\begin{tikzpicture}
\begin{axis}[
  width=13.5cm, height=7cm,
  xlabel={Variable $x$},
  ylabel={Variable $y$},
  xmin=0, xmax=9.5, ymin=-2, ymax=18,
  xtick={0,2,4,6,8},
  ytick={0,4,8,12,16},
  grid=both, grid style={gray!25},
  axis lines=left,
  tick label style={font=\small},
  label style={font=\small},
  legend style={at={(0.02,0.98)}, anchor=north west, font=\small, fill=white, draw=gray!40},
]
\addplot[only marks, mark=*, mark size=2.8pt, color=black!80] 
  coordinates {(1,2)(2,4)(3,6)(4,8)(5,10)(6,12)(7,2)(8,16)};
\addplot[thick, color=biocolor, domain=0:9, samples=2] {2*x};
\addplot[thick, dashed, color=bioaccent, domain=0:9, samples=2] {1.286*x + 1.714};
\legend{,Sans point aberrant : $y = 2x$ ($r=1$),Avec point aberrant : $y=1{,}29x+1{,}71$ ($r\approx 0{,}66$)}
\end{axis}
\end{tikzpicture}
\end{center}
\textbf{Constat} : un \textbf{seul point} aberrant (le point situé très en dessous de la tendance) suffit à :
\begin{itemize}[leftmargin=*,itemsep=1pt]
  \item « aplatir » la pente ($a$ passe de $2$ à $1{,}29$) ;
  \item dégrader fortement le coefficient de corrélation ($r$ passe de $1$ à $0{,}66$) ;
  \item fausser toute prévision faite avec cette droite.
\end{itemize}
\textbf{Leçon} : avant tout calcul, \textbf{repérer visuellement les points aberrants}. Vérifier s'il s'agit d'une erreur de saisie ou d'une valeur biologique exceptionnelle à traiter séparément.\\
\textbf{Reflexe} : tracer systématiquement le nuage et comparer la droite « avec » et « sans » les points suspects.
\end{vignetteAlerte}

\newpage

% =========================================================
% ANNEXE D bis – TABLEAU COMPARATIF DES GRAPHIQUES
% =========================================================
\section*{\color{biocolor}Annexe D bis -- Tableau comparatif des graphiques}
\emph{Synthèse : quel graphique choisir, pour quel type de données, avec quels avantages et limites.}

\vspace{0.3cm}

\begin{tcolorbox}[enhanced, colback=biolight, colframe=biocolor,
                  title={\faChartBar\ \ \textbf{Choisir le bon graphique}},
                  coltitle=white, fonttitle=\bfseries\large,
                  boxrule=0.8pt, arc=4pt, left=8pt, right=8pt, top=8pt, bottom=8pt,
                  drop fuzzy shadow=black!15]
\small
\renewcommand{\arraystretch}{1.35}
\begin{tabular}{|p{2.4cm}|p{2.6cm}|p{2.6cm}|p{3.4cm}|p{3.2cm}|}
\hline
\rowcolor{biocolor!20}
\textbf{Graphique} & \textbf{Type de variable} & \textbf{Ce qu'il montre} & \textbf{Intérêt principal} & \textbf{Limites / précautions} \\
\hline
\textbf{Bâtons} (D1) & Qualitative ou quantitative discrète & Effectif / fréquence de chaque modalité & Compare rapidement les catégories ; repère le mode & Peu adapté aux variables continues \\
\hline
\textbf{Circulaire} (D2) & Qualitative (peu de modalités) & Part de chaque catégorie dans le total & Vision « composition » très parlante pour un public non spécialiste & Difficile à lire au-delà de 6--7 modalités ; compare mal deux séries \\
\hline
\textbf{Histogramme} (D3) & Quantitative continue (regroupée en classes) & Forme de la distribution (symétrie, mode, outliers) & Révèle la structure globale ; oriente le choix moyenne/médiane & Sensible au choix des classes (amplitude, bornes) \\
\hline
\textbf{Courbe des FCC} (D4) & Quantitative (discrète ou continue) & Cumul des fréquences ; lecture directe de $M_e$, $Q_1$, $Q_3$ & Répond à « quel \% est inférieur à $x$ ? » ; complément de l'histogramme & Moins visuelle seule ; ne montre pas la forme \\
\hline
\textbf{Boîte à moustaches} (D5) & Quantitative (une ou plusieurs séries) & 5 indicateurs : min, $Q_1$, $M_e$, $Q_3$, max ; outliers & Compare plusieurs séries en un coup d'œil ; robuste & Ne montre pas la forme exacte (bimodalité invisible) \\
\hline
\textbf{Nuage de points + droite} (D6) & Deux variables quantitatives ($x$, $y$) & Liaison entre deux variables ; tendance linéaire & Prédiction, étalonnage ; détection d'outliers et de non-linéarité & Nécessite de tracer le nuage \emph{avant} de calculer $r$ ; $r$ ne mesure que la linéarité \\
\hline
\end{tabular}
\end{tcolorbox}

\vspace{0.3cm}

\begin{vignetteCours}{D10 -- Quel graphique pour quel objectif ? (arbre de décision simplifié)}
\begin{enumerate}[leftmargin=*,itemsep=2pt]
  \item \textbf{Une seule variable, qualitative ou discrète ?} \textrightarrow\ Diagramme en bâtons (ou circulaire si peu de modalités).
  \item \textbf{Une seule variable, quantitative continue ?}
    \begin{itemize}[leftmargin=*,itemsep=0pt]
      \item Pour voir la \emph{forme} de la distribution \textrightarrow\ Histogramme.
      \item Pour lire $M_e$, $Q_1$, $Q_3$ sans calcul \textrightarrow\ Courbe des FCC.
      \item Pour \emph{comparer plusieurs séries} \textrightarrow\ Boîtes à moustaches juxtaposées.
    \end{itemize}
  \item \textbf{Deux variables quantitatives ?} \textrightarrow\ Nuage de points, puis, si la tendance est linéaire, droite de régression.
  \item \textbf{Toujours vérifier} : tracer d'abord, calculer ensuite. Un graphique mal choisi peut masquer une information essentielle.
\end{enumerate}
\end{vignetteCours}

\newpage

% =========================================================
% ANNEXE E – VIGNETTES "GRILLES DE CALCUL"
% =========================================================
\section*{\color{bioaccent}Annexe E -- Vignettes « grilles de calcul semi-manuel »}
\emph{Grilles prêtes à remplir, dimensionnées selon les données exactes de la compétition. Les formules à appliquer sont rappelées dans les vignettes méthodes de l'Annexe B.}

\vspace{0.3cm}

% ---------- E1 ----------
\begin{vignetteMethode}{E1 -- Grille $\bar{x}$, $V$, $\sigma$ (8 classes de l'épreuve 1)}
\emph{Nombre de lignes utiles : \textbf{8} (une par classe $[12\,;13[$ à $[19\,;20[$).}
\emph{Formules : voir la vignette \textbf{M2}.}

\vspace{0.2cm}
\renewcommand{\arraystretch}{1.5}
\begin{center}
\small
\begin{tabular}{|c|c|c|c|c|c|c|c|}
\hline
$i$ & Classe & $x_i$ & $n_i$ & $n_i x_i$ & $x_i^2$ & $n_i x_i^2$ & FCC \\
\hline
1 & $[12\,;13[$ & 12{,}5 & & & & & \\ \hline
2 & $[13\,;14[$ & 13{,}5 & & & & & \\ \hline
3 & $[14\,;15[$ & 14{,}5 & & & & & \\ \hline
4 & $[15\,;16[$ & 15{,}5 & & & & & \\ \hline
5 & $[16\,;17[$ & 16{,}5 & & & & & \\ \hline
6 & $[17\,;18[$ & 17{,}5 & & & & & \\ \hline
7 & $[18\,;19[$ & 18{,}5 & & & & & \\ \hline
8 & $[19\,;20[$ & 19{,}5 & & & & & \\ \hline
\multicolumn{2}{|r|}{\textbf{Total}} & --- & $N=40$ & $\sum n_i x_i =$ & --- & $\sum n_i x_i^2 =$ & 1 \\ \hline
\end{tabular}
\end{center}
\end{vignetteMethode}

% ---------- E2 ----------
\begin{vignetteMethode}{E2 -- Grille sur les 40 valeurs brutes (pré-transcrites)}
\emph{Nombre de lignes utiles : \textbf{40} (les 40 valeurs du contexte commun, déjà pré-remplies pour éviter toute erreur de recopie).}
\emph{Formules : voir la vignette \textbf{M2}.}

\vspace{0.2cm}
\renewcommand{\arraystretch}{1.25}
\begin{center}
\footnotesize
\begin{tabular}{|c|c|c|@{\hspace{0.4cm}}|c|c|c|}
\hline
$i$ & $x_i$ & $x_i^2$ & $i$ & $x_i$ & $x_i^2$ \\
\hline
1  & 12{,}1 & & 21 & 15{,}9 & \\ \hline
2  & 12{,}6 & & 22 & 16{,}0 & \\ \hline
3  & 13{,}0 & & 23 & 16{,}1 & \\ \hline
4  & 13{,}3 & & 24 & 16{,}3 & \\ \hline
5  & 13{,}6 & & 25 & 16{,}4 & \\ \hline
6  & 13{,}9 & & 26 & 16{,}5 & \\ \hline
7  & 14{,}0 & & 27 & 16{,}6 & \\ \hline
8  & 14{,}2 & & 28 & 16{,}8 & \\ \hline
9  & 14{,}4 & & 29 & 16{,}9 & \\ \hline
10 & 14{,}5 & & 30 & 17{,}0 & \\ \hline
11 & 14{,}7 & & 31 & 17{,}2 & \\ \hline
12 & 14{,}9 & & 32 & 17{,}4 & \\ \hline
13 & 15{,}0 & & 33 & 17{,}5 & \\ \hline
14 & 15{,}1 & & 34 & 17{,}7 & \\ \hline
15 & 15{,}2 & & 35 & 17{,}9 & \\ \hline
16 & 15{,}3 & & 36 & 18{,}1 & \\ \hline
17 & 15{,}4 & & 37 & 18{,}4 & \\ \hline
18 & 15{,}5 & & 38 & 18{,}6 & \\ \hline
19 & 15{,}6 & & 39 & 18{,}9 & \\ \hline
20 & 15{,}8 & & 40 & 19{,}1 & \\ \hline
\multicolumn{3}{|c|}{$\sum_{1}^{20} x_i = \quad \sum_{1}^{20} x_i^2 =$} 
  & \multicolumn{3}{c|}{$\sum_{21}^{40} x_i = \quad \sum_{21}^{40} x_i^2 =$} \\ \hline
\end{tabular}
\end{center}
\end{vignetteMethode}

% ---------- E3 ----------
\begin{vignetteMethode}{E3 -- Grille bivariée $\bar{x}$, $\bar{y}$, $V(x)$, $V(y)$, $\mathrm{Cov}(x,y)$, $r$ (6 points de l'épreuve 5)}
\emph{Nombre de lignes utiles : \textbf{6} (les 6 couples $(x_i\,;y_i)$ de l'étalonnage).}
\emph{Formules : voir la vignette \textbf{M5}.}

\vspace{0.2cm}
\renewcommand{\arraystretch}{1.5}
\begin{center}
\small
\begin{tabular}{|c|c|c|c|c|c|}
\hline
$i$ & $x_i$ & $y_i$ & $x_i^2$ & $y_i^2$ & $x_i y_i$ \\
\hline
1 & 0 & 0{,}05 & & & \\ \hline
2 & 2 & 0{,}22 & & & \\ \hline
3 & 4 & 0{,}38 & & & \\ \hline
4 & 6 & 0{,}57 & & & \\ \hline
5 & 8 & 0{,}71 & & & \\ \hline
6 & 10 & 0{,}90 & & & \\ \hline
\multicolumn{1}{|r|}{\textbf{Total}} & $\sum x_i =$ & $\sum y_i =$ & $\sum x_i^2 =$ & $\sum y_i^2 =$ & $\sum x_i y_i =$ \\ \hline
\end{tabular}
\end{center}
\end{vignetteMethode}

% ---------- E4 ----------
\begin{vignetteMethode}{E4 -- Grille d'interpolation pour $Q_1$, $M_e$, $Q_3$ (8 classes)}
\emph{Nombre de lignes utiles : \textbf{3} (les 3 indicateurs à déterminer).}
\emph{Méthode : voir la vignette \textbf{M3}.}

\vspace{0.2cm}
\renewcommand{\arraystretch}{1.7}
\begin{center}
\small
\begin{tabular}{|c|c|c|c|c|c|c|}
\hline
\textbf{Indicateur} & \textbf{Niveau} & \textbf{Classe} & $a$ & $b$ & $F_{i-1}$ & $f_i$ \\
\hline
$Q_1$ & 0{,}25 & & & & & \\[0.15cm] \hline
$M_e$ & 0{,}50 & & & & & \\[0.15cm] \hline
$Q_3$ & 0{,}75 & & & & & \\[0.15cm] \hline
\end{tabular}
\end{center}
\end{vignetteMethode}

\newpage

% ---------- E5 (enrichie) ----------
\begin{vignetteMethode}{E5 -- Récapitulatif des sommes et résultats}
\emph{Nombre de lignes utiles : \textbf{17} (une par quantité à calculer).}
\emph{Pour chaque quantité, le rôle, l'intérêt et l'interprétation sont précisés.}

\vspace{0.2cm}
\renewcommand{\arraystretch}{1.45}
\begin{center}
\small
\begin{tabular}{|p{3cm}|p{9cm}|p{3cm}|}
\hline
\textbf{Quantité} & \textbf{Rôle, intérêt et interprétation} & \textbf{Valeur calculée} \\ \hline
Effectif total $N$ & 
\tiny\textbf{Rôle : } base de tous les calculs (moyennes, variances). 
Intérêt : vérifier la cohérence des données. 
Interprétation : nombre total d'observations (ici $N=40$). & \\ \hline
$\sum x_i$ & 
\tiny\textbf{Rôle : } nécessaire pour calculer $\bar{x}$. 
Intérêt : permet de vérifier les calculs. 
Interprétation : somme brute des concentrations. & \\ \hline
$\sum x_i^2$ & 
\tiny\textbf{Rôle : } utilisée dans la formule rapide de la variance. 
Intérêt : évite de calculer les écarts à la moyenne. 
Interprétation : grandeur liée à la dispersion et à la moyenne quadratique. & \\ \hline
$\bar{x}$ & 
\tiny\textbf{Rôle : } indicateur de position centrale. 
Intérêt : résume la tendance centrale de $x$. 
Interprétation : concentration moyenne en glucose. & \\ \hline
$V(x)$ & 
\tiny\textbf{Rôle : } mesure la dispersion autour de la moyenne. 
Intérêt : quantifie l'hétérogénéité des valeurs. 
Interprétation : plus $V(x)$ est grande, plus les valeurs sont dispersées. & \\ \hline
$\sigma_x = \sqrt{V(x)}$ & 
\tiny\textbf{Rôle : } mesure de dispersion dans la même unité que $x$. 
Intérêt : interprétable directement. 
Interprétation : écart moyen à la moyenne. & \\ \hline
$\sum y_i$ & 
\tiny\textbf{Rôle : } pour calculer $\bar{y}$. 
Intérêt : calcul de la moyenne des absorbances. 
Interprétation : somme brute des absorbances. & \\ \hline
$\sum y_i^2$ & 
\tiny\textbf{Rôle : } pour calculer $V(y)$. 
Intérêt : évite les écarts à la moyenne. 
Interprétation : lié à la dispersion des absorbances. & \\ \hline
$\sum x_i y_i$ & 
\tiny\textbf{Rôle : } pour calculer la covariance. 
Intérêt : mesure la liaison entre $x$ et $y$. 
Interprétation : si positive, $x$ et $y$ varient dans le même sens. & \\ \hline
$\bar{y}$ & 
\tiny\textbf{Rôle : } position centrale de $y$. 
Intérêt : nécessaire pour la droite de régression. 
Interprétation : absorbance moyenne. & \\ \hline
$V(y)$ & 
\tiny\textbf{Rôle : } dispersion de $y$. 
Intérêt : calcul de $\sigma_y$ et de $r$. 
Interprétation : variabilité des absorbances. & \\ \hline
$\sigma_y = \sqrt{V(y)}$ & 
\tiny\textbf{Rôle : } dispersion de $y$ dans son unité. 
Intérêt : pour le calcul de $r$. 
Interprétation : écart moyen des absorbances à la moyenne. & \\ \hline
$\mathrm{Cov}(x,y)$ & 
\tiny\textbf{Rôle : } mesure la variation conjointe de $x$ et $y$. 
Intérêt : signe indique le sens de la relation, valeur pour calculer $a$ et $r$. 
Interprétation : si $>0$, relation positive ; si $<0$, négative ; si proche de 0, pas de relation linéaire. & \\ \hline
Coefficient $a$ (pente) & 
\tiny\textbf{Rôle : } quantifie la variation de $y$ quand $x$ augmente de 1 unité. 
Intérêt : permet de prédire $y$ à partir de $x$. 
Interprétation : augmentation d'absorbance par g/L de glucose. & \\ \hline
Coefficient $b$ (ordonnée) & 
\tiny\textbf{Rôle : } valeur de $y$ quand $x=0$. 
Intérêt : ajustement de la droite. 
Interprétation : absorbance théorique pour concentration nulle (souvent proche de 0, mais peut être décalage). & \\ \hline
Coefficient $r$ & 
\tiny\textbf{Rôle : } mesure la force de la liaison linéaire entre $x$ et $y$. 
Intérêt : évaluer la qualité de l'ajustement linéaire. 
Interprétation : entre $-1$ et $1$ ; $|r|$ proche de $1$ indique une forte corrélation linéaire ; proche de $0$, pas de corrélation linéaire. & \\ \hline
Équation de la droite & 
\tiny\textbf{Rôle : } modèle de prédiction. 
Intérêt : estimer $y$ pour un $x$ donné ou inversement. 
Interprétation : outil d'étalonnage. & $y = $ \\ \hline
\end{tabular}
\end{center}
\end{vignetteMethode}

\newpage

% =========================================================
% ANNEXE F – LEXIQUE (tableau sécable sur plusieurs pages)
% =========================================================
\section*{\color{biocolor}Annexe F -- Lexique des termes statistiques}
\emph{Tableau à double colonne : terme technique et définition.}
\emph{Les termes sont classés par ordre alphabétique. Le tableau se poursuit sur la page suivante si nécessaire.}

\vspace{0.3cm}

{\small
\renewcommand{\arraystretch}{1.3}
\begin{longtable}{|p{3cm}|p{12.5cm}|}
\hline
\textbf{Terme} & \textbf{Définition} \\
\hline
\endfirsthead
\hline
\textbf{Terme} & \textbf{Définition} \\
\hline
\endhead
\hline
\multicolumn{2}{|r|}{\emph{Suite du lexique page suivante\dots}} \\
\endfoot
\hline
\endlastfoot
Amplitude (d'une classe) & Différence entre la borne supérieure et la borne inférieure d'une classe : $b-a$. Constante si les classes sont régulières. \\
\hline
Boîte à moustaches & Représentation graphique qui résume une série par 5 valeurs : $\min$, $Q_1$, $M_e$, $Q_3$, $\max$. \\
\hline
Caractère & Ce que l'on mesure sur un individu : quantitatif (numérique) ou qualitatif (catégorie). \\
\hline
Centre de classe & Milieu d'une classe $[a\,;b[$ : $x_i = (a+b)/2$. Utilisé comme représentant de la classe dans les calculs. \\
\hline
Classe & Intervalle $[a\,;b[$ regroupant les valeurs d'une série continue ou nombreuse. \\
\hline
Classe modale & Classe (dans une série regroupée) qui a le plus grand effectif. \\
\hline
Coefficient de corrélation linéaire ($r$) & Nombre entre $-1$ et $1$ qui mesure la force et le sens de la liaison \emph{linéaire} entre deux variables. $|r|$ proche de $1$ = forte corrélation ; proche de $0$ = pas de corrélation linéaire. \\
\hline
Covariance & Mesure de la variation conjointe de deux variables : $\mathrm{Cov}(x,y)=\frac{1}{N}\sum x_i y_i - \bar{x}\bar{y}$. Son signe indique le sens de la relation. \\
\hline
Diagramme circulaire & Représentation en secteurs dont l'angle est proportionnel à la fréquence : $\theta_i = f_i \times 360^\circ$. Utilisé pour un caractère qualitatif. \\
\hline
Diagramme en bâtons & Représentation par barres verticales pour une variable discrète ou qualitative. \\
\hline
Distribution & Manière dont les valeurs d'une série se répartissent (forme, centre, dispersion). \\
\hline
Écart interquartile ($I_Q$) & Différence entre le troisième et le premier quartile : $I_Q = Q_3 - Q_1$. Mesure la dispersion des 50\,\% centraux. Robuste aux valeurs extrêmes. \\
\hline
Écart-type ($\sigma$) & Racine carrée de la variance : $\sigma = \sqrt{V}$. Exprimé dans la même unité que la variable. Mesure la dispersion moyenne autour de la moyenne. \\
\hline
Effectif ($n_i$) & Nombre d'individus d'une modalité ou d'une classe. \\
\hline
Effectif total ($N$) & Somme de tous les effectifs : $N = \sum n_i$. \\
\hline
Étendue & Différence entre le maximum et le minimum : $\max - \min$. \\
\hline
Fréquence ($f_i$) & Proportion d'individus dans une modalité : $f_i = n_i / N$. La somme des fréquences vaut 1. \\
\hline
Fréquence cumulée croissante (FCC) & Somme des fréquences jusqu'à une modalité donnée : $F_i = f_1 + f_2 + \dots + f_i$. La dernière FCC vaut toujours 1. \\
\hline
Histogramme & Représentation graphique pour une variable continue : rectangles jointifs dont l'\emph{aire} est proportionnelle à l'effectif. \\
\hline
Individu & Élément de la population étudiée (ici, un échantillon sérique). \\
\hline
Interpolation linéaire & Méthode pour estimer une valeur (ex. $Q_1$) à l'intérieur d'une classe à partir de la courbe des FCC ou de la formule $x = a + (b-a)\frac{p - F_{i-1}}{f_i}$. \\
\hline
Médiane ($M_e$) & Valeur qui partage la série ordonnée en deux moitiés : 50\,\% en dessous, 50\,\% au-dessus. Robuste aux valeurs extrêmes. \\
\hline
Mode & Valeur (ou classe) de plus grand effectif. Une série peut avoir plusieurs modes. \\
\hline
Moindres carrés & Méthode qui détermine la droite $y = ax + b$ minimisant la somme des carrés des écarts verticaux entre points et droite. \\
\hline
Moyenne ($\bar{x}$) & Somme de toutes les valeurs divisée par l'effectif total : $\bar{x} = \frac{1}{N}\sum n_i x_i$. Sensible aux valeurs extrêmes. \\
\hline
Non paramétrique & Se dit d'une méthode qui ne suppose pas de forme particulière pour la distribution (ex. médiane, quartiles). \\
\hline
Nuage de points & Représentation graphique des couples $(x_i\,;y_i)$ dans le plan. \\
\hline
Outlier (valeur aberrante) & Valeur très éloignée des autres, qui peut fausser la moyenne, l'écart-type et la corrélation. À repérer sur un graphique avant tout calcul. \\
\hline
Paramétrique & Se dit d'une méthode qui suppose une forme de distribution (souvent normale) et utilise moyenne et écart-type. \\
\hline
Population & Ensemble complet des individus étudiés. \\
\hline
Premier quartile ($Q_1$) & Valeur telle que 25\,\% des données lui sont inférieures ou égales. \\
\hline
Quartiles & Les trois valeurs $Q_1$, $M_e$, $Q_3$ qui partagent la série ordonnée en quatre parties égales (25\,\% chacune). \\
\hline
Régression linéaire & Ajustement d'une droite $y = ax + b$ à un nuage de points par la méthode des moindres carrés. \\
\hline
Robuste & Se dit d'un indicateur peu sensible aux valeurs extrêmes (ex. médiane, écart interquartile). \\
\hline
Série statistique & Liste des valeurs observées d'un caractère sur une population. \\
\hline
Symétrie (d'une distribution) & Propriété d'une distribution dont le côté gauche est l'image miroir du côté droit. Se lit sur la boîte à moustaches ou l'histogramme. \\
\hline
Troisième quartile ($Q_3$) & Valeur telle que 75\,\% des données lui sont inférieures ou égales. \\
\hline
Valeur aberrante & Voir \emph{outlier}. \\
\hline
Variance ($V$) & Moyenne des carrés des écarts à la moyenne : $V = \frac{1}{N}\sum n_i(x_i-\bar{x})^2 = \frac{1}{N}\sum n_i x_i^2 - \bar{x}^2$. \\
\hline
\end{longtable}
}

\vspace{0.3cm}

\begin{vignetteCours}{F1 -- Repères mnémotechniques}
\begin{itemize}[leftmargin=*,itemsep=1pt]
  \item \textbf{Moyenne et écart-type} : couple « sensible » ; utile si distribution symétrique sans outlier.
  \item \textbf{Médiane et écart interquartile} : couple « robuste » ; à préférer si distribution asymétrique ou avec outliers.
  \item \textbf{Paramétrique} = suppose une loi (normale) ; \textbf{non paramétrique} = ne suppose rien.
  \item \textbf{Outlier} = valeur aberrante ; à repérer \emph{avant} tout calcul.
  \item \textbf{Corrélation ($r$)} mesure un lien \emph{linéaire} uniquement : $r \approx 0$ n'exclut pas un lien non linéaire (ex. parabole).
  \item \textbf{Moindres carrés} = méthode pour trouver la « meilleure » droite $y = ax + b$.
\end{itemize}
\end{vignetteCours}

\end{document}